% GENERATED FILE. Edit canonical lessons, then run npm run generate:books.
% canonical-source-hash: fnv1a64:ebf160e53295d5fb
% canonical-lessons: GU-C99-pelu, GU-C99-avu, GU-C99-evu, GU-C99-darek, GU-C99-koi

\chapter{That, Like this, Like that, Every, Someone}
\label{ch:gu-that-like-this-like-that-every-someone}
\hlchaptermodality{99}

\begin{chapteropening}
\textbf{By the end of this chapter:} \emph{I can say that, like this, like that, every and someone.}

Complete the last lesson of chapter 99: I can say that, like this, like that, every and someone.
\end{chapteropening}

\section[\texorpdfstring{\gu{પેલું} (pelũ)}{પેલું (pelũ)}]{\gu{પેલું} (pelũ) — that (over there)}

\label{lesson:GU-C99-pelu}



\begin{quote}
Before the new one: say the Gujarati for Bhadarvo, then the Gujarati for Aso.
\end{quote}

\subsection*{You'll want to know: \gu{પેલું}}
\textbf{\gu{પેલું}} (\emph{pelũ}) — \textquotedblleft{}that (over there)\textquotedblright{}.

\begin{cousinweb}
\textbf{\gu{પેલું}} (\emph{pelũ}) is a tadbhava. CDIAL lists it under a reconstructed \emph{*pratilla}, \textquotedblleft{}opposite\textquotedblright{}, built on Sanskrit \emph{prati}, \textquotedblleft{}against\textquotedblright{} (CDIAL 8594), and again under \emph{pratikūla}, \textquotedblleft{}contrary, opposite\textquotedblright{} (CDIAL 8545).
\end{cousinweb}

\begin{grammarlens}[title={What you've built}]
One more word for this chapter.
\end{grammarlens}

\subsection*{Your turn}
\begin{itemize}
\raggedright
  \item \emph{Say it:} \emph{pelũ}
  \item \emph{Say it:} \emph{pelũ}, once more
  \item \emph{Say it:} say \emph{pelũ}
  \item \emph{Recall:} say the Gujarati for Bhadarvo, then the Gujarati for Aso, then say \emph{pelũ} again
\end{itemize}

\begin{tcolorbox}[breakable,colback=teal!4,colframe=teal!35!black,title={Before you move on}]
What does \gu{પેલું} mean? (\textquotedblleft{}That (over there)\textquotedblright{}.) Say it once more.
\end{tcolorbox}
\section[\texorpdfstring{\gu{આવું} (āvũ)}{આવું (āvũ)}]{\gu{આવું} (āvũ) — like this, such}

\label{lesson:GU-C99-avu}



\begin{quote}
Before the new one: say the Gujarati for Aso, then the Gujarati for that.
\end{quote}

\subsection*{You'll want to know: \gu{આવું}}
\textbf{\gu{આવું}} (\emph{āvũ}) — \textquotedblleft{}like this, such\textquotedblright{}.

\begin{cousinweb}
The origin of \textbf{\gu{આવું}} (\emph{āvũ}) is not settled. CDIAL, searched by form and by meaning, has no entry for it.
\end{cousinweb}

\begin{grammarlens}[title={What you've built}]
One more word for this chapter.
\end{grammarlens}

\subsection*{Your turn}
\begin{itemize}
\raggedright
  \item \emph{Say it:} \emph{āvũ}
  \item \emph{Say it:} \emph{āvũ}, once more
  \item \emph{Say it:} say \emph{āvũ}
  \item \emph{Recall:} say the Gujarati for Aso, then the Gujarati for that, then say \emph{āvũ} again
\end{itemize}

\begin{tcolorbox}[breakable,colback=teal!4,colframe=teal!35!black,title={Before you move on}]
What does \gu{આવું} mean? (\textquotedblleft{}Like this, such\textquotedblright{}.) Say it once more.
\end{tcolorbox}
\section[\texorpdfstring{\gu{એવું} (evũ)}{એવું (evũ)}]{\gu{એવું} (evũ) — like that, such}

\label{lesson:GU-C99-evu}



\begin{quote}
Before the new one: say the Gujarati for that, then the Gujarati for like this.
\end{quote}

\subsection*{You'll want to know: \gu{એવું}}
\textbf{\gu{એવું}} (\emph{evũ}) — \textquotedblleft{}like that, such\textquotedblright{}.

\begin{cousinweb}
The origin of \textbf{\gu{એવું}} (\emph{evũ}) is not settled. CDIAL, searched by form and by meaning, has no entry for it.
\end{cousinweb}

\begin{grammarlens}[title={What you've built}]
One more word for this chapter.
\end{grammarlens}

\subsection*{Your turn}
\begin{itemize}
\raggedright
  \item \emph{Say it:} \emph{evũ}
  \item \emph{Say it:} \emph{evũ}, once more
  \item \emph{Say it:} say \emph{evũ}
  \item \emph{Recall:} say the Gujarati for that, then the Gujarati for like this, then say \emph{evũ} again
\end{itemize}

\begin{tcolorbox}[breakable,colback=teal!4,colframe=teal!35!black,title={Before you move on}]
What does \gu{એવું} mean? (\textquotedblleft{}Like that, such\textquotedblright{}.) Say it once more.
\end{tcolorbox}
\section[\texorpdfstring{\gu{દરેક} (darek)}{દરેક (darek)}]{\gu{દરેક} (darek) — every, each}

\label{lesson:GU-C99-darek}



\begin{quote}
Before the new one: say the Gujarati for like this, then the Gujarati for like that.
\end{quote}

\subsection*{You'll want to know: \gu{દરેક}}
\textbf{\gu{દરેક}} (\emph{darek}) — \textquotedblleft{}every, each\textquotedblright{}.

\begin{cousinweb}
The origin of \textbf{\gu{દરેક}} (\emph{darek}) is not settled. CDIAL, searched by form and by meaning, has no entry for it.
\end{cousinweb}

\begin{grammarlens}[title={What you've built}]
One more word for this chapter.
\end{grammarlens}

\subsection*{Your turn}
\begin{itemize}
\raggedright
  \item \emph{Say it:} \emph{darek}
  \item \emph{Say it:} \emph{darek}, once more
  \item \emph{Say it:} say \emph{darek}
  \item \emph{Recall:} say the Gujarati for like this, then the Gujarati for like that, then say \emph{darek} again
\end{itemize}

\begin{tcolorbox}[breakable,colback=teal!4,colframe=teal!35!black,title={Before you move on}]
What does \gu{દરેક} mean? (\textquotedblleft{}Every, each\textquotedblright{}.) Say it once more.
\end{tcolorbox}
\section[\texorpdfstring{\gu{કોઈ} (koī)}{કોઈ (koī)}]{\gu{કોઈ} (koī) — someone, anyone}

\label{lesson:GU-C99-koi}



\begin{quote}
Before the new one: say the Gujarati for like that, then the Gujarati for every.
\end{quote}

\subsection*{You'll want to know: \gu{કોઈ}}
\textbf{\gu{કોઈ}} (\emph{koī}) — \textquotedblleft{}someone, anyone\textquotedblright{}.

\begin{cousinweb}
\textbf{\gu{કોઈ}} (\emph{koī}) is a tadbhava, a word inherited from Sanskrit and reshaped by sound change. It comes from \emph{kaścid}, \textquotedblleft{}someone, anyone\textquotedblright{} (CDIAL 2967).
\end{cousinweb}

\begin{grammarlens}[title={What you've built}]
One more word for this chapter.
\end{grammarlens}

\subsection*{Your turn}
\begin{itemize}
\raggedright
  \item \emph{Say it:} \emph{koī}
  \item \emph{Say it:} \emph{koī}, once more
  \item \emph{Say it:} say \emph{koī}
  \item \emph{Recall:} say the Gujarati for like that, then the Gujarati for every, then say \emph{koī} again
\end{itemize}

\begin{tcolorbox}[breakable,colback=teal!4,colframe=teal!35!black,title={Before you move on}]
What does \gu{કોઈ} mean? (\textquotedblleft{}Someone, anyone\textquotedblright{}.) Say it once more.
\end{tcolorbox}
